\documentclass[a4paper,12pt]{article}




%\abovedisplayshortskip=-12pt plus 3pt
%\setlength{\abovedisplayskip}{-5pt}
%\setlength{\belowdisplayskip}{-5pt}
%\linespread{1.05}
%\voffset=-14cm
%\hoffset=-0.5cm

\usepackage{bm, amsmath, amsfonts, amssymb, url}
\usepackage{titlesec, titletoc}
\usepackage{graphicx,subfigure}
\usepackage{appendix, indentfirst, makeidx, array}
\usepackage{booktabs}
\usepackage{fancyhdr}
\pagestyle{fancy}
\usepackage{fullpage} 

\addtolength{\headsep}{+1cm}
\addtolength{\oddsidemargin}{+0.5cm}
\addtolength{\evensidemargin}{-0.5cm}
%\usepackage[cam,a4,center]{crop}
%\usepackage{geometry}


%\linespread{1.3}

%\fancypagestyle{plain}{
%\fancyhf{}
%\renewcommand{\headrulewidth}{0pt}
%%\fancyfoot[c]{\thepage}
%}


%\addtocontents{toc}{\protect\thispagestyle{empty}}

%%%-------page style ------------

%\setlength{\baselineskip}{20pt plus0.3pt minus0.2pt}
%\setlength{\marginparsep}{0pt}
%\setlength{\marginparwidth}{1.5cm}

%\setlength{\oddsidemargin}{3cm}
%\setlength{\evensidemargin}{3cm}
%%%-------------------------------


\renewcommand{\labelenumi}{(\alph{enumi})}
\renewcommand{\arraystretch}{1.2}
\titleformat{\chapter}[display]{\Large\bf}{\LARGE{Chapter~\thechapter}}{2mm}{}
\titleformat{\section}{\large\bf}{\thesection}{1.5mm}{}
\titleformat{\subsection}{\bfseries\itshape}{\thesubsection}{1.5mm}{}
%self-define symbol
\newcommand{\proposition}[1]{\bigskip\noindent\fbox{\begin{minipage}{0.95\textwidth}~\begin{minipage}{0.95\textwidth}#1\end{minipage}\\\smallskip\end{minipage}}\bigskip\vline\vline\vline}

\newcommand{\excercise}{\addtocounter{section}{+1}
\subsection*{\underline{Ex~\thesection}}}
\newcommand{\stag}[1]{\tag{\arabic{chapter}.\arabic{equation}#1}}

\newcommand{\qdo}[0]{\quad}
\newcommand{\qdt}[0]{\quad\quad}
\newcommand{\qdthr}[0]{\quad\quad\quad}
\newcommand{\qdf}[0]{\quad\quad\quad\quad}
\newcommand{\qdfiv}[0]{\quad\quad\quad\quad\quad}

\newcommand{\norm}[1]{\|#1\|}
%\newcommand{\sdelta}{{\scriptstyle\Delta}}
\newcommand{\sdelta}{\Delta}
\newcommand{\bra}[1]{\langle#1|}
\newcommand{\ket}[1]{|#1\rangle }
\newcommand{\bket}[1]{\Big |#1\Big\rangle }
\newcommand{\braket}[2]{\langle#1|#2\rangle }
\newcommand{\proj}[2]{|#1\rangle\langle#2|}
\newcommand{\expect}[2]{\langle #1\rangle_{#2}}
\newcommand{\Kappa}[0]{\bm{\mathcal{K}}}
\newcommand{\parity}[0]{\bm{\mathcal{P}}}
\newcommand{\parityl}[0]{\bm{p}}
\newcommand{\charge}[0]{\bm{\mathcal{C}}}


\makeindex
%-----------------------------------%

\begin{document}

\abovedisplayskip=-6pt plus 3pt minus 0pt
\belowdisplayskip=8pt plus 3pt minus 0pt

\abovedisplayshortskip=-6pt plus 3pt minus 0pt
\belowdisplayshortskip=8pt plus 3pt minus 0pt

\title{Problem Set 4 (Due: Oct. 19,2011)}
\maketitle

Your homework should be submitted to TA (Chih-Han Lin) at the  beginning of the course. Generally, You have a week to finish the problem set since it was revealed. If you cannot submit it on time, you can still bring your homework to TA's office (L325, Institute of Atomic and Molecular Sciences, in NTU campus) and put it into TA's mailbox which is on the top of the shoe cabinet in L325. Always remember to remind TA to check your homework by e-mail (\url{clin@ltl.iams.sinica.edu.tw}) if you don't submit it in class.

%\abovedisplayskip=0pt %plus 3pt minus 3pt
%\belowdisplayskip=2pt %plus 3pt minus 3pt
%
%\abovedisplayshortskip=0pt %plus 3pt minus 3pt
%\belowdisplayshortskip=2pt %plus 3pt minus 3pt

\excercise
Show that delta function $\delta(x)$ can be expressed as

\begin{enumerate}
\item $$
\delta(x)=\frac{1}{\pi}\lim_{N\rightarrow \infty}\frac{\sin Nx}{x}.
$$

\item $$
\delta(x)=\frac{1}{2}\frac{d^2}{dx^2}|x|.
$$
 
\end{enumerate}


\excercise

Show that 

$$
\bra{x}\mathbf{P}^2\ket{x'}=\left(\frac{\hbar}{i}\right)^2\delta''(x-x'),
$$
 
\noindent and hence generalize to 

$$
\bra{x}F(\mathbf{P})\ket{x'}=F\left(\frac{\hbar}{i}\frac{d}{dx}\right)\delta(x-x').
$$



\excercise

Consider a quantum system with a state $\ket{\psi}$ that is translated by a distance $\xi$ by the unitary operator $\mathbf{U}(\mathbf{P}; \xi)=\exp(-i\xi\mathbf{P}/\hbar)$, i.e.

$$
\ket{\psi'}=\mathbf{U}(\mathbf{P}; \xi)\ket{\psi}=e^{-\frac{i}{\hbar}\xi\mathbf{P}}\ket{\psi}.
$$

Find the q-representation of $\ket{\psi'}$, which is the wave function $\psi'(x)=\braket{x}{\psi'}=\bra{x}\mathbf{U}(\mathbf{P}; \xi)\ket{\psi}$, explicitly and interpret the result.




\end{document}
